%==============================================================================
%  wavemaker_sponge_BC.tex
%
%  Drop-in subsection for main.tex, intended to be \input just after
%  Section 8 ("Gridap Implementation"), i.e. after the framed global residual
%  \eqref{eq: multifield general residual}.  It documents the three practical
%  mechanisms the solver adds on top of the bare weak form: (i) the treatment
%  of horizontal boundary conditions, (ii) the internal wavemaker mass source,
%  and (iii) the sponge (relaxation) absorbing layers.
%
%  It relies on the macros already defined in the preamble of main.tex
%  (\Oh, \nablah, \uj, \vi, \Ns, \sumj, \sumi, \VmM, \VMij, \Phij, \ddt, ...).
%  If compiled standalone, copy those \newcommand definitions across.
%==============================================================================

\subsection*{Wave Generation, Absorption, and Boundary Conditions}

The global residual \eqref{eq: multifield general residual} derived above describes wave
\emph{propagation} in the interior of $\Oh$, but a usable simulation additionally requires a way to
\emph{inject} energy into the domain, to \emph{absorb} it before it reaches the artificial mesh
edges, and to prescribe what happens \emph{on} those edges. In the solver these are handled by three
mechanisms, each of which augments the weak form with one explicit, closed-form term. Crucially, all
three are \emph{linear} in the unknowns, so they leave the Newton solve structure of
Section~8 unchanged.

\subsubsection*{Boundary conditions}

Two boundary integrals were generated and then discarded during the weak-form derivation: the
continuity flux $\oint_{\partial\Oh} q H\uj\cdot\mathbf{n}\,d\Gamma$ arising from the divergence
integration by parts \eqref{eq: semi-discrete mass continuity equation weak form}, and the analogous
momentum flux $\nabla_h\cdot(\dots)$ from the horizontal pressure-gradient term (Section~6). The
solver realises exactly the conditions under which those integrals vanish:

\begin{itemize}
    \item \textbf{Solid walls (essential / Dirichlet).} On a rigid lateral wall the no-penetration
    condition $\uj\cdot\mathbf{n}=0$ holds for \emph{every} vertical layer $j$. For a wall normal to
    $y$ this means the essential constraint
    \begin{equation*}
        u_j^y\big|_{\partial\Oh^{\text{wall}}} = 0, \qquad j = 0,1,\dots,\Ns,
    \end{equation*}
    imposed directly on the corresponding components of the \texttt{MultiFieldFESpace} (and
    symmetrically $u_j^x=0$ for a wall normal to $x$ when a fully closed basin is required). Both
    boundary integrals then vanish identically, and the free-surface elevation $\eta$ (equivalently
    $H$) is left \emph{unconstrained} --- it satisfies a natural (Neumann) condition, as is physically
    correct for a moving free surface against a vertical wall.

    \textbf{Implementation note.} On a Cartesian mesh the wall is tagged by its edge \emph{and its two
    end corners}; omitting the corner tags leaves the corner degrees of freedom unconstrained and
    triggers an exponential instability. The solid-wall constraint therefore uses the full tag set
    (four corners $+$ the two edge interiors) per wall.

    \item \textbf{Open boundaries (natural).} At an open end of a flume no essential condition is
    prescribed; the boundary flux is left as the natural (homogeneous) term. Physical radiation of
    the outgoing waves through such an open edge is \emph{not} enforced by the boundary condition
    itself but by placing a sponge layer (below) in front of it, so that the wave amplitude reaching
    the edge is already negligible and the residual boundary reflection is harmless.
\end{itemize}

\subsubsection*{Internal wavemaker: a mass source in the continuity equation}

Rather than prescribing an inflow boundary condition, waves are generated \emph{internally} by adding
a localised mass source $S(\mathbf{x},t)$ to the right-hand side of the depth-integrated mass
continuity equation \eqref{eq: semi-discrete mass continuity residual}:
\begin{equation}
    \ddt{H} + \sumj\nablah\cdot\left(H\uj\right)\Phij = S(\mathbf{x},t).
    \label{eq: continuity with source}
\end{equation}
Weighting by the trial function $q$ and integrating over $\Oh$ exactly as in
\eqref{eq: semi-discrete mass continuity equation weak form}, the source contributes a single extra
term to the global residual:
\begin{equation}
    \boxed{\; R_{\text{wm}} = -\int_{\Oh} q\, S(\mathbf{x},t)\, d\Oh \;}
    \label{eq: wavemaker residual contribution}
\end{equation}
The source is a time-harmonic Gaussian pulse in space,
\begin{equation}
    S(\mathbf{x},t) = 2\,A\,\omega\; \mathcal{G}(\mathbf{x})\;\cos(\omega t),
    \qquad \omega = \frac{2\pi}{T},
    \label{eq: wavemaker source profile}
\end{equation}
with target wave amplitude $A$ and angular frequency $\omega$ set by the wave period $T$. The spatial
footprint $\mathcal{G}(\mathbf{x})$ selects the type of generated wavefield:
\begin{subequations}\label{eq: wavemaker footprints}
    \begin{align}
        \text{Line source (long-crested / plane waves):} && \mathcal{G}(\mathbf{x}) &= \exp\!\left[-\left(\frac{x-x_{\text{wm}}}{\sigma_{\text{wm}}}\right)^{2}\right], \label{eq: wavemaker line}\\[1mm]
        \text{Point source (radial / ring waves):} && \mathcal{G}(\mathbf{x}) &= \exp\!\left[-\frac{(x-x_{\text{wm}})^{2}+(y-y_{\text{wm}})^{2}}{\sigma_{\text{wm}}^{2}}\right]. \label{eq: wavemaker point}
    \end{align}
\end{subequations}
Here $(x_{\text{wm}}, y_{\text{wm}})$ is the generator location, and $\sigma_{\text{wm}}$ is the
Gaussian half-width controlling how sharply the source is concentrated (a narrow footprint,
$\sigma_{\text{wm}}\sim1.5$~m by default, approximates a clean single-frequency paddle while
remaining smooth enough to be resolved by the horizontal mesh). The line source
\eqref{eq: wavemaker line} has no $y$-dependence and therefore emits straight wave crests; the point
source \eqref{eq: wavemaker point} emits circular fronts.

\paragraph{The factor of 2.} Because the source injects mass symmetrically, it radiates waves in
\emph{both} directions away from the generator line/point. Only half of the emitted energy travels
toward the region of interest; the leading factor $2$ in \eqref{eq: wavemaker source profile}
compensates for this split so that the wave reaching the working area has the prescribed amplitude
$A$. The counter-propagating half is absorbed by the sponge layer on the opposite side.

\paragraph{Solver implementation.} The profile \eqref{eq: wavemaker source profile} is built by the
factory routines \texttt{make\_wavemaker\_line}$(x_{\text{wm}},A,T,k;\sigma_{\text{wm}})$ and
\texttt{make\_wavemaker\_point}$(x_{\text{wm}},y_{\text{wm}},A,T;\sigma_{\text{wm}})$, which return a
closure $S(\mathbf{x},t)$. This closure is stored on the problem struct (\texttt{prob.wm\_src}) and,
at each time $t$, wrapped into a Gridap \texttt{CellField} evaluated on the integration mesh; the term
\eqref{eq: wavemaker residual contribution} is then appended to the continuity residual.

\subsubsection*{Sponge (relaxation) layers: absorption in the momentum equation}

To prevent spurious reflections from the artificial domain edges, absorbing \emph{sponge} layers add a
Rayleigh-type linear damping to the horizontal momentum equation. For vertical layer $i$, the damping
force acts through the same vertical mass matrix $\VMij$ that appears in the acceleration term
\eqref{eq: horizontal momentum semi-discrete acceleration residual}, ensuring that all vertical modes
are damped consistently:
\begin{equation}
    \mathbf{Sp}_i^V \equiv \mu(\mathbf{x})\sumj \VMij\,\uj .
    \label{eq: sponge semi-discrete term}
\end{equation}
Weighting by $\vi$ and summing over $i$ as in Section~8, and collecting the layers into the
multi-field vectors $\mathbf{U}$, $\mathbf{V}$ of
\eqref{eq: def U and V horizontal velocity multifield FE spaces}, the sponge contributes
\begin{equation}
    \boxed{\; R_{\text{sp}} = \sumi\left[\int_{\Oh}\mu(\mathbf{x})\left(\sumj\VMij\,\uj\right)\vi\;d\Oh\right] = \int_{\Oh}\mu(\mathbf{x})\left(\VmM\,\mathbf{U}\right)\cdot\mathbf{V}\;d\Oh \;}
    \label{eq: sponge residual contribution}
\end{equation}
The damping coefficient $\mu(\mathbf{x})\ge0$ is nonzero only inside the boundary layers and grows
\emph{quadratically} from zero at the inner edge of each layer to its maximum $\mu_{\max}$ at the
domain boundary. With up to four layers (left, right, bottom, top) of widths $w_L, w_R, w_B, w_T$ on a
rectangular domain $[x_0,x_1]\times[y_0,y_1]$, the profile is the pointwise maximum of the individual
edge ramps:
\begin{equation}
    \mu(\mathbf{x}) = \max\Big\{\, \mu_L, \mu_R, \mu_B, \mu_T \,\Big\},
    \label{eq: sponge profile max}
\end{equation}
with, for example, the left and bottom contributions
\begin{align}
    \mu_L(\mathbf{x}) &= \mu_{\max}\left(\frac{x_0 + w_L - x}{w_L}\right)^{2}
        \quad\text{for } x < x_0 + w_L, \qquad 0 \text{ otherwise}, \label{eq: sponge left}\\[1mm]
    \mu_B(\mathbf{x}) &= \mu_{\max}\left(\frac{y_0 + w_B - y}{w_B}\right)^{2}
        \quad\text{for } y < y_0 + w_B, \qquad 0 \text{ otherwise}, \label{eq: sponge bottom}
\end{align}
and symmetrically for the right ($x > x_1 - w_R$) and top ($y > y_1 - w_T$) edges. The argument of
each square is the normalised penetration depth into the layer, so $\mu$ ramps smoothly from $0$
(inner edge, no damping, no reflection) to $\mu_{\max}$ (outer wall, full damping). Taking the
\emph{maximum} rather than the sum in \eqref{eq: sponge profile max} clamps the overlapping corner
regions at $\mu_{\max}$ instead of $2\mu_{\max}$, avoiding an over-damped corner that would itself
reflect.

\paragraph{Design remarks.}
\begin{itemize}
    \item The sponge damps \emph{momentum} only; it does not appear in the continuity equation, so it
    removes kinetic energy while leaving mass conservation intact away from the generator.
    \item Because \eqref{eq: sponge residual contribution} is linear in $\mathbf{U}$, it adds directly
    and exactly to the momentum Jacobian block (no linearisation needed) and does not affect the
    effective mass matrix.
    \item The required strength scales with the wave regime: deep-water ($kd$ large) cases need larger
    $\mu_{\max}$ or wider layers than shallow-water cases, since a weak sponge lets amplitude build up
    over long runs.
\end{itemize}

\paragraph{Solver implementation.} The profile \eqref{eq: sponge profile max} is built by
\texttt{make\_sponge\_2D}$(\text{domain}, w_L, w_R, w_B, w_T, \mu_{\max})$, which returns the closure
$\mu(\mathbf{x})$. It is stored on the problem struct (\texttt{prob.mu\_sponge}) and, being
time-independent, is wrapped once into a static \texttt{CellField}; the term
\eqref{eq: sponge residual contribution} is appended to each momentum mode of the residual (and to the
Jacobian).

\subsubsection*{Augmented global residual}

Collecting \eqref{eq: wavemaker residual contribution} and
\eqref{eq: sponge residual contribution}, the complete residual passed to Gridap is the interior form
\eqref{eq: multifield general residual} together with the source and sponge contributions:
\begin{framed}
    \begin{multline}
        \text{Global Residual}_{\text{full}} = \text{Global Residual}\big|_{\eqref{eq: multifield general residual}}
        \;-\; \int_{\Oh} q\, S(\mathbf{x},t)\, d\Oh
        \;+\; \int_{\Oh} \mu(\mathbf{x})\left(\VmM\,\mathbf{U}\right)\cdot\mathbf{V}\; d\Oh .
        \label{eq: augmented global residual}
    \end{multline}
\end{framed}
The wavemaker term contributes only to the load vector (right-hand side), being independent of the
unknowns; the sponge term contributes symmetrically to both residual and Jacobian. Neither disturbs
the tensor-product separation of Section~7--8: the wavemaker is a purely horizontal forcing, and the
sponge factors as a horizontal scalar field $\mu(\mathbf{x})$ times the vertical mass matrix $\VmM$.
